\documentclass{article}
\usepackage{lmodern}
\usepackage{amsmath}
\usepackage{listings}
\usepackage{fullpage}
\usepackage{parskip}
\usepackage{xcolor}
\lstset{
	basicstyle=\ttfamily,
	columns=fullflexible,
	frame=single,
	breaklines=true,
	postbreak=\mbox{\textcolor{red}{$\hookrightarrow$}\space},
}
\begin{document}
\title{Experiment `p\_4\_2\_a\_min\_max\_seq' Results}
\author{\today}
\date{}
\maketitle
\textbf{Experiment outcome: }SUCCESS
\\\textbf{Bad responses: }0
\\\textbf{Responses containing}~\texttt{assume}~\textbf{: }0
\\\textbf{Responses containing}~\texttt{expect}~\textbf{: }0
\\\textbf{Resolution attempts: }1
\\\textbf{Hard fails (resolution): }0
\\\textbf{Soft fails (resolution): }0
\\\textbf{Verification attempts: }1
\section*{Problem Specification}
\textbf{Problem name: }p\_4\_2\_a\_min\_max\_seq
\\\textbf{Natural language statement: }Write a method that takes a sequence of integers and returns the smallest and largest of the inputs.
\\\textbf{Method signature: }p\_4\_2\_a\_min\_max\_seq(inputs: seq<int>) returns (min: int, max: int)
\subsection*{Ensures}
\begin{itemize}
\item \texttt{forall i :: 0 <= i < |inputs| ==> min <= inputs[i]}
\item \texttt{forall i :: 0 <= i < |inputs| ==> max >= inputs[i]}
\item \texttt{exists i :: 0 <= i < |inputs| \&\& min == inputs[i]}
\item \texttt{exists i :: 0 <= i < |inputs| \&\& max == inputs[i]}
\end{itemize}
\subsection*{Requires}
\begin{itemize}
\item \texttt{0 < |inputs|}
\end{itemize}
\clearpage
\section*{GenAI interactions}
Below you will find all interactions between the `user' (program) and the `assistant' (OpenAI).
\subsection*{Program $\rightarrow$ GenAI}
\begin{lstlisting}
You are given the following task to perform in Dafny:

Write a method that takes a sequence of integers and returns the smallest and largest of the inputs.

The signature should be:

p_4_2_a_min_max_seq(inputs: seq<int>) returns (min: int, max: int)

The method should respect the following contract:

requires 0 < |inputs|, ensures forall i :: 0 <= i < |inputs| ==> min <= inputs[i], ensures forall i :: 0 <= i < |inputs| ==> max >= inputs[i], ensures exists i :: 0 <= i < |inputs| && min == inputs[i], ensures exists i :: 0 <= i < |inputs| && max == inputs[i]

Produce and show only the Dafny body of this method, including the curly braces that surround it. Do not show the signature nor contract. You must not use 'assume' nor 'expect' anywhere in your code.
\end{lstlisting}
\subsection*{GenAI $\rightarrow$ Program}
\textbf{System fingerprint: }null
\\\textbf{ID: }chatcmpl-Dme4pWxueOvzeC3iHYBG7QaT8afNw
\\\textbf{Model: }gpt-5.5-2026-04-23
\\\textbf{Created at: }1780486171
\\\textbf{Finish reason: }stop
\\\textbf{Usage: }422 tokens in, and 1613 tokens out
\begin{lstlisting}
{
  min := inputs[0];
  max := inputs[0];

  ghost var minIndex := 0;
  ghost var maxIndex := 0;
  var i := 1;

  while i < |inputs|
    invariant 1 <= i <= |inputs|
    invariant 0 <= minIndex < i
    invariant 0 <= maxIndex < i
    invariant min == inputs[minIndex]
    invariant max == inputs[maxIndex]
    invariant forall j :: 0 <= j < i ==> min <= inputs[j]
    invariant forall j :: 0 <= j < i ==> max >= inputs[j]
    decreases |inputs| - i
  {
    if inputs[i] < min {
      min := inputs[i];
      minIndex := i;
    }

    if inputs[i] > max {
      max := inputs[i];
      maxIndex := i;
    }

    i := i + 1;
  }
}
\end{lstlisting}
\section*{Final Program}
\begin{lstlisting}
method p_4_2_a_min_max_seq(inputs: seq<int>) returns (min: int, max: int)
	requires 0 < |inputs|
	ensures forall i :: 0 <= i < |inputs| ==> min <= inputs[i]
	ensures forall i :: 0 <= i < |inputs| ==> max >= inputs[i]
	ensures exists i :: 0 <= i < |inputs| && min == inputs[i]
	ensures exists i :: 0 <= i < |inputs| && max == inputs[i]
{
  min := inputs[0];
  max := inputs[0];

  ghost var minIndex := 0;
  ghost var maxIndex := 0;
  var i := 1;

  while i < |inputs|
    invariant 1 <= i <= |inputs|
    invariant 0 <= minIndex < i
    invariant 0 <= maxIndex < i
    invariant min == inputs[minIndex]
    invariant max == inputs[maxIndex]
    invariant forall j :: 0 <= j < i ==> min <= inputs[j]
    invariant forall j :: 0 <= j < i ==> max >= inputs[j]
    decreases |inputs| - i
  {
    if inputs[i] < min {
      min := inputs[i];
      minIndex := i;
    }

    if inputs[i] > max {
      max := inputs[i];
      maxIndex := i;
    }

    i := i + 1;
  }
}
\end{lstlisting}
\section*{Total Token Usage}
\textbf{Input tokens: }422
\\\textbf{Output tokens: }1613
\\\textbf{Reasoning tokens: }1361
\\\textbf{Sum of `total tokens': }2035
\section*{Experiment Timings}
\textbf{Overall Experiment} started at 1780486170881, ended at 1780486210996, lasting 40115ms (40.12 seconds)
\\\textbf{Iteration \#1} started at 1780486170888, ended at 1780486210996, lasting 40108ms (40.11 seconds)
\end{document}
